\title{\textbf{Simulación de colas: servidor CI/CD M/M/1 y SOC con\\ prioridades y buffer finito}}
\author{Joshtin Mejías (31071016), Angel Bandres (29524378) y Andrés Crespo (31259533)\\[0.2em]
\small Proyecto de Simulación de Colas --- Prof. Lisbeth Torres}
\date{24 de noviembre de 2026}
\maketitle
\thispagestyle{empty}

\begin{abstract}
\noindent
Este informe documenta la construcción y validación de un simulador de eventos
discretos para dos sistemas de líneas de espera. El primero es un servidor
CI/CD modelado como M/M/1 con $\lambda=0{,}25$ trabajos/min, $\mu=1/3$ y
$\rho=0{,}75$. El segundo es un Centro de Operaciones de Seguridad (SOC) con
$s=3$ analistas, alertas críticas y preventivas, prioridad no expropiativa,
servicios no exponenciales y un buffer finito de $K=8$ espacios. La simulación
del M/M/1 reproduce los valores analíticos dentro de los intervalos de confianza
del 95\,\% y satisface la Ley de Little. El SOC base presenta una pérdida de
alertas de $0{,}385$\,\% y una utilización conjunta de $68{,}87$\,\%. El
Escenario A, con un aumento del $40$\,\% en alertas preventivas, eleva la
pérdida a $2{,}220$\,\% ($+477{,}0$\,\% relativo) sin saturar teóricamente a los
analistas, pero con un fuerte aumento de la congestión. En el Escenario B,
ampliar el buffer a $K=16$ y contratar un cuarto analista reducen la pérdida de
forma similar bajo la demanda base; sin embargo, el cuarto analista reduce
drásticamente los tiempos de espera y domina también bajo sobrecarga. Se
recomienda contratar el cuarto analista manteniendo $K=8$.
\end{abstract}
